\section{Literatür verisiyle yazılım uygulamaları}
\label{sec:spss-b13}

\textbf{Amaç ve veri.} \texttt{b13.csv} ile birinci dönem notu
\texttt{g1} ve yıl sonu notu \texttt{g3} arasındaki doğrusal ilişkiyi
inceleyin \citep{cortez2008data}. Pearson korelasyonunda
$H_0:\rho=0$, sabit terimli basit regresyonda $H_0:\beta_1=0$
sınanır. Alternatifler iki yönlüdür; anlamlılık düzeyi 0,05,
katsayıların güven düzeyi yüzde 95'tir. Açıklayıcı değişken G1,
yanıt değişkeni G3'tür; bu yön bir nedensellik iddiası değildir.

\textbf{Ortak veri.} Karşılaştırmanın veri dosyası
\nolinkurl{companion/spss/csv/b13.csv} olup 395 satır ve 10 sütun
içerir. Kayıtların tamamı kullanılır; sıfır yıl sonu notu bulunan
38 kayıt eksik sayılmaz. B06'nın seçim göstergeleri ve ağırlığı
kullanılmaz. Sayısal uyum Excel içe aktarmanın doğrulandığı anlamına
gelmez. Kodları \texttt{companion} klasörünü içeren paket kökünden
çalıştırın; veri sözlüğü için bkz. sayfa~\pageref{sec:spss-guide}.

\subsection{IBM SPSS uygulaması}
\textbf{Başlamadan önce.} Aşağıdaki veri açma bloğunu
\texttt{EXECUTE} satırı dahil çalıştırın; sonra menü veya analiz
komutlarıyla devam edin. Filtre, ağırlık ve dosya bölme kapalı
olmalıdır. Dosya açılamazsa \texttt{/FILE} yolunu düzeltip yeniden
açın; önceki bölümden kalan etkin veriyle devam etmeyin.

\begin{spssadimlar}
\spssadim{Önce \emph{Graphs $\to$ Legacy Dialogs $\to$ Scatter/Dot} yolunda basit saçılım grafiği oluşturun. Yatay eksene \texttt{g1}, dikey eksene \texttt{g3} koyun; doğrusal örüntüyü ve aykırı kayıtları inceleyin.}
\spssadim{\emph{Analyze $\to$ Correlate $\to$ Bivariate} yolunu açın. \texttt{g1} ve \texttt{g3} değişkenlerini \emph{Variables} alanına taşıyın.}
\spssadim{Doğrusal ilişkiyi inceleyen bu uygulamada \emph{Pearson} ve \emph{Two-tailed} seçin. \emph{OK} ile tabloyu alın; yöntem seçimini çıkan $p$ değerine göre değiştirmeyin.}
\spssadim{Regresyon için \emph{Analyze $\to$ Regression $\to$ Linear} yolunu açın. \emph{Dependent} alanına \texttt{g3}, açıklayıcı değişken alanına \texttt{g1} koyun; yöntem \emph{Enter} olsun.}
\spssadim{\emph{Statistics} içinde \emph{Estimates}, \emph{Model fit} ve yüzde 95 \emph{Confidence intervals} seçin; \emph{Continue} seçin.}
\spssadim{\emph{Plots} içinde Y alanına \texttt{*ZRESID}, X alanına \texttt{*ZPRED} yerleştirin; standartlaştırılmış artık histogramını da seçin.}
\spssadim{\emph{Continue $\to$ OK} ile modeli çalıştırın. \emph{Coefficients} tablosunun standartlaştırılmamış B sütununu ve \emph{Model Summary} tablosunun $R^2$ değerini okuyun.}
\end{spssadimlar}

{\raggedright
\noindent\textbf{Komutların görevi.} \texttt{CORRELATIONS} Pearson ilişkisini, \texttt{REGRESSION} yıl sonu notunun doğrusal modelini hesaplar. \texttt{SCATTERPLOT} ve \texttt{RESIDUALS} model artıklarını incelemek içindir.\par}

\par\noindent\begin{minipage}{\linewidth}
\textbf{Uygulama kodu:} \texttt{b13}. Veri açma ve tanımlama:

\begin{lstlisting}[style=prismSPSS]
SET DECIMAL=DOT.
GET DATA /TYPE=TXT
 /FILE='companion/spss/csv/b13.csv'
 /ENCODING='UTF8'
 /ARRANGEMENT=DELIMITED /DELCASE=LINE
 /DELIMITERS="," /QUALIFIER='"'
 /FIRSTCASE=2
 /VARIABLES=id F8.0 school F8.0 sex F8.0 age F8.0
 studytime F8.0 absences F8.0 g1 F8.0 g2 F8.0
 g3 F8.0 pass10 F8.0.
FILTER OFF.
WEIGHT OFF.
SPLIT FILE OFF.
VARIABLE LABELS
 g1 'Ilk donem matematik notu'
 g3 'Yil sonu matematik notu'.
VARIABLE LEVEL
 id school sex pass10 (NOMINAL)
 studytime (ORDINAL)
 age absences g1 g2 g3 (SCALE).
EXECUTE.
\end{lstlisting}
\end{minipage}\par

\par\noindent\begin{minipage}{\linewidth}
\textbf{Korelasyon, model ve grafikler.} Veri açıldıktan sonra çalıştırın.

\begin{lstlisting}[style=prismSPSS]
DESCRIPTIVES VARIABLES=g1 g3
 /STATISTICS=MEAN STDDEV MIN MAX.
GRAPH /SCATTERPLOT(BIVAR)=g1 WITH g3.
CORRELATIONS /VARIABLES=g1 g3
 /PRINT=TWOTAIL /MISSING=PAIRWISE.
REGRESSION /STATISTICS=COEFF R ANOVA CI(95)
 /DEPENDENT=g3 /METHOD=ENTER g1
 /SCATTERPLOT=(*ZRESID,*ZPRED)
 /RESIDUALS=HISTOGRAM(ZRESID).
\end{lstlisting}

\end{minipage}\par

\begin{outputguidebox}
\textbf{Paylaşılan SPSS tablolarıyla doğrulanan değerler.}
\emph{Descriptive Statistics} tablosunda iki değişken için de
$n=395$ vardır. G1 ortalaması 10,91, standart sapması 3,319;
G3 ortalaması 10,42, standart sapması 4,581'dir.
\emph{Correlations} tablosunda $r=0{,}801$ ve iki yönlü
$p<0{,}001$ görülür. \emph{Model Summary}, $R^2=0{,}642$,
düzeltilmiş $R^2=0{,}641$ ve tahminin standart hatası 2,743 verir.

\emph{ANOVA} tablosunda $F(1,393)=705{,}842$ ve $p<0{,}001$;
\emph{Coefficients} tablosunda sabit $-1{,}653$, G1 eğimi
1,106 ve standartlaştırılmış Beta 0,801'dir. Eğim aralığı
$[1{,}024;\ 1{,}188]$ olarak gösterilir. B eğimi ile Beta
aynı sütun değildir. SPSS'in $p<0{,}001$ gösterimi sıfır demek
değildir; R ve Python daha küçük değerleri bilimsel gösterimle verir.

\emph{Residuals Statistics} tablosunda ham artık standart sapması
2,740, standartlaştırılmış artık standart sapması 0,999'dur.
Bunlar modelin 2,743 standart hatasıyla karıştırılmaz: artık
standart sapmasının böleni 394, model hata varyansının böleni
393'tür. Burada ZRESID, ham artığın model standart hatasına
bölünmesidir; standart sapmasının tam 1 olması beklenmez.
\end{outputguidebox}


\par\noindent\begin{minipage}{\linewidth}
\subsection{R uygulaması}
\label{sec:r-gercek-b13}

\texttt{lm} sabit terimli modeli kurar. Aşağıdaki blokları sırayla
çalıştırın; tahmin ve artıklar ham kayıtlardan ayrı nesnelerdir.

\begin{lstlisting}[style=prismR]
veri <- read.csv("companion/spss/csv/b13.csv")
stopifnot(nrow(veri) == 395, ncol(veri) == 10,
  !anyNA(veri), !anyDuplicated(veri$id),
  all(veri$pass10 == as.integer(veri$g3 >= 10)))
korelasyon <- cor.test(veri$g1, veri$g3,
  method="pearson", alternative="two.sided")
model <- lm(g3 ~ g1, data=veri)
ozet <- summary(model)
tahmin <- fitted(model)
artik <- residuals(model)
zpred <- as.numeric(scale(tahmin))
zresid <- artik/sigma(model)
print(c(satir=nrow(veri), sutun=ncol(veri),
  eksik=sum(is.na(veri)), sifir_not=sum(veri$g3 == 0)))
\end{lstlisting}
\end{minipage}\par

\textbf{Dosya yolu sorunu yaşarsanız.} İlk satır yerine aşağıdaki
komutla \texttt{b13.csv} dosyasını seçin; ardından kontrol bloğundan
devam edin.
\begin{lstlisting}[style=prismR]
veri <- read.csv(file.choose())
\end{lstlisting}

\par\noindent\begin{minipage}{\linewidth}
\textbf{Modelin sayısal özeti.} Küçük $p$ değerlerini sıfıra
yuvarlamamak için 15 anlamlı basamakla yazdırın.

\begin{lstlisting}[style=prismR]
print(c(r=unname(korelasyon$estimate),
  p_korelasyon=korelasyon$p.value,
  R2=ozet$r.squared, duzeltilmis_R2=ozet$adj.r.squared,
  model_sh=sigma(model),
  F=unname(ozet$fstatistic["value"]),
  df_model=unname(ozet$fstatistic["numdf"]),
  df_artik=df.residual(model),
  beta=coef(model)[["g1"]]*sd(veri$g1)/sd(veri$g3),
  G1_10_tahmin=unname(predict(model, data.frame(g1=10)))
), digits=15)
print(cbind(ozet$coefficients,
            confint(model, level=.95)), digits=15)
print(anova(model), digits=15)
\end{lstlisting}

\textbf{Çıktıyı okuma.} Paylaşılan R ekranlarında
$r=0{,}801467932017414$, $R^2=0{,}642350846052270$ ve
G1=10 için tahmin 9,40975711892797 bulunmuştur. Katsayı tablosu
hem sabit hem eğim için standart hata, $t$, $p$ ve güven aralığı
verir. R burada yeniden çalıştırılmamış; paylaşılan çıktılar incelenmiştir.
\end{minipage}\par

\par\noindent\begin{minipage}{\linewidth}
\textbf{Betimsel ve artık özetleri.} SPSS'in artık tablosuyla
karşılaştırma için aynı ölçeklemeyi koruyun.

\begin{lstlisting}[style=prismR]
betimsel <- function(degerler) {
  c(n=length(degerler), minimum=min(degerler),
    maksimum=max(degerler), ortalama=mean(degerler),
    ss=sd(degerler))
}
print(rbind(G1=betimsel(veri$g1), G3=betimsel(veri$g3)),
      digits=15)
print(rbind(tahmin=betimsel(tahmin), artik=betimsel(artik),
  zpred=betimsel(zpred), zresid=betimsel(zresid)), digits=15)
\end{lstlisting}

\textbf{Çıktıyı okuma.} R'de ZRESID aralığı yaklaşık
$[-4{,}236510;\ 1{,}828148]$, standart sapması 0,998730'dur.
Artık ortalamalarının $10^{-16}$ civarında görünmesi sayısal
hesap hassasiyetiyle ilgilidir; anlamlı bir ortalama sapması değildir.
\end{minipage}\par

\par\noindent\begin{minipage}{\linewidth}
\textbf{Grafikleri üretme.} Paylaşılan R görseli üç paneli de içerir.

\begin{lstlisting}[style=prismR]
par(mfrow=c(1, 3))
plot(veri$g1, veri$g3, xlab="G1", ylab="G3",
     main="G1 - G3", pch=16)
abline(model, col="blue", lwd=2)
plot(zpred, zresid, xlab="ZPRED", ylab="ZRESID",
     main="Artik - Tahmin", pch=16)
abline(h=0, lty=2)
hist(zresid, breaks=15, main="Artik histogrami",
     xlab="ZRESID")
par(mfrow=c(1, 1))
\end{lstlisting}
\end{minipage}\par

\par\noindent\begin{minipage}{\linewidth}
\subsection{Python uygulaması}
\label{sec:python-b13}

\texttt{linregress} eğim, sabit ve eğimin standart hatasını verir.
Eğim aralığında $n-2=393$ serbestlik derecesi kullanılır.
Blokları sırayla çalıştırın; \texttt{ddof=1} örneklem standart
sapmasını seçer.

\begin{lstlisting}[style=prismPythonApp]
import pandas as pd
import numpy as np
from scipy import stats

veri = pd.read_csv("companion/spss/csv/b13.csv")
assert veri.shape == (395, 10)
assert not veri.isna().any().any()
assert veri.id.is_unique
assert veri.pass10.eq((veri.g3 >= 10).astype(int)).all()
ilk_not = veri.g1.to_numpy()
son_not = veri.g3.to_numpy()
korelasyon = stats.pearsonr(ilk_not, son_not)
model = stats.linregress(ilk_not, son_not)
serbestlik = len(veri)-2
kritik = stats.t.ppf(.975, serbestlik)
egim_araligi = (model.slope
    + np.array([-1, 1])*kritik*model.stderr)
tahmin = model.intercept+model.slope*ilk_not
artik = son_not-tahmin
model_sh = np.sqrt(np.sum(artik**2)/serbestlik)
zresid = artik/model_sh
zpred = (tahmin-tahmin.mean())/tahmin.std(ddof=1)
\end{lstlisting}
\end{minipage}\par

\par\noindent\begin{minipage}{\linewidth}
\textbf{Sayısal çıktı.} Etiketler paylaşılan terminal çıktısıyla aynıdır.

\begin{lstlisting}[style=prismPythonApp]
print("Satir:", len(veri), "Sutun:", veri.shape[1])
print("Eksik:", int(veri.isna().sum().sum()))
print("Sifir not:", int((veri.g3 == 0).sum()))
ozet = {
    "Pearson r": korelasyon.statistic,
    "Korelasyon p": korelasyon.pvalue,
    "Sabit": model.intercept, "Egim": model.slope,
    "Egim standart hatasi": model.stderr,
    "Egim p": model.pvalue,
    "Egim GA alt": egim_araligi[0],
    "Egim GA ust": egim_araligi[1],
    "R2": model.rvalue**2,
    "Duzeltilmis R2": 1-(1-model.rvalue**2)
        * (len(veri)-1)/serbestlik,
    "Model standart hatasi": model_sh,
    "F": (model.slope/model.stderr)**2,
    "Artik serbestlik derecesi": serbestlik,
    "G1=10 ortalama tahmini": model.intercept+10*model.slope}
for ad, deger in ozet.items():
    print(f"{ad}: {deger:.15g}")
\end{lstlisting}

\textbf{Çıktıyı okuma.} Python ekranında 395 satır, 10 sütun,
sıfır eksik değer ve 38 sıfır not vardır. Korelasyon ve eğim için
$p\approx9{,}00143\times10^{-90}$ bulunmuştur. \texttt{.15g},
15 ondalık basamak değil 15 anlamlı basamak ister; çok küçük
değerleri bilimsel gösterimle korur.
\end{minipage}\par

\par\noindent\begin{minipage}{\linewidth}
\textbf{Grafikleri üretme.} Çizgi G1'e göre sıralanarak çizilir.

\begin{lstlisting}[style=prismPythonApp]
import matplotlib.pyplot as plt

sekil, eksenler = plt.subplots(1, 3, figsize=(15, 4))
sirali = np.argsort(ilk_not)
eksenler[0].scatter(ilk_not, son_not, alpha=.6)
eksenler[0].plot(ilk_not[sirali], tahmin[sirali],
                 color="red")
eksenler[0].set(xlabel="G1", ylabel="G3", title="G1 - G3")
eksenler[1].scatter(zpred, zresid, alpha=.6)
eksenler[1].axhline(0, color="black", linestyle="--")
eksenler[1].set(xlabel="ZPRED", ylabel="ZRESID",
                title="Artik - Tahmin")
eksenler[2].hist(zresid, bins=15, edgecolor="black")
eksenler[2].set(xlabel="ZRESID", ylabel="Frekans",
                title="Artik histogrami")
plt.tight_layout(); plt.show()
\end{lstlisting}
\end{minipage}\par

\Needspace{12\baselineskip}
\subsection{Sonuçların karşılaştırılması ve ortak rapor}
13 Eylül 2026'da paylaşılan SPSS tabloları ve grafikleri, R ekranları
ve Python terminal çıktısı ile üç panelli Python grafiği incelenmiştir.
Python'un 14 sayısal model özeti R'deki karşılıklarıyla uyumludur:
$p$ dışındaki değerlerde $10^{-10}$ mutlak, çok küçük $p$ değerlerinde
$10^{-10}$ bağıl tolerans kullanılmıştır. Bütün basamaklarda eşitlik
iddia edilmez. Python kodu proje CSV'siyle ayrıca çalıştırılmış;
SPSS sonuçları gösterilen yuvarlama hassasiyetlerinde uyumlu bulunmuştur.
Sabitin aralığı ve ayrıntılı ANOVA gibi Python terminalinde ayrıca
yazdırılmayan değerler için SPSS ve R çıktıları esas alınmıştır.
Aşağıdaki tablolar bu çıktıların kitap düzenine uyarlanmış özetidir;
ekran görüntüsü değildir.

\begin{table}[H]
\centering
\caption{B13'te korelasyon ve basit regresyonun uyumu}
\label{tab:b13-dogrulanmis-uyum}
\begin{tabular}{@{}lr@{}}
\toprule
Ölçü & Değer \\
\midrule
Gözlem sayısı & 395 \\
Pearson $r$ & 0,801468 \\
$R^2$ & 0,642351 \\
Düzeltilmiş $R^2$ & 0,641441 \\
Modelin artık standart hatası & 2,743359 \\
$F(1,393)$ & 705,842247 \\
Korelasyon ve eğim için $p$ (yaklaşık) & $9{,}00143\times10^{-90}$ \\
Standartlaştırılmış Beta & 0,801468 \\
G1=10 için ortalama tahmini & 9,409757 \\
\bottomrule
\end{tabular}
\par\smallskip
{\footnotesize Sabit terimli, tek açıklayıcılı bu modelde
$R^2=r^2$, standartlaştırılmış Beta $=r$ ve eğimin $t$ istatistiğinin
karesi modelin $F$ istatistiğidir. Korelasyon ve eğim testleri
bağımsız iki kanıt olarak sayılmaz.\par}
\end{table}

\begin{table}[H]
\centering
\caption{B13'te standartlaştırılmamış regresyon katsayıları}
\label{tab:b13-dogrulanmis-katsayi}
\begin{tabular}{@{}lrr@{}}
\toprule
Ölçü & Sabit & G1 eğimi \\
\midrule
B & $-1{,}652804$ & 1,106256 \\
Standart hata & 0,474745 & 0,041639 \\
$t(393)$ & $-3{,}481452$ & 26,567692 \\
İki yönlü $p$ (yaklaşık) & 0,000555 & $9{,}00143\times10^{-90}$ \\
\%95 güven aralığı: alt & $-2{,}586162$ & 1,024393 \\
\%95 güven aralığı: üst & $-0{,}719445$ & 1,188119 \\
\bottomrule
\end{tabular}
\par\smallskip
{\footnotesize Güven aralıkları klasik regresyon modeline dayanır;
dayanıklı standart hatalar kullanılmamıştır. Katsayıların yorumunda
artık grafikleri ve bağımsızlık varsayımı ayrıca değerlendirilir.\par}
\end{table}

\begin{table}[H]
\centering
\caption{B13'te regresyon varyans çözümlemesi}
\label{tab:b13-dogrulanmis-anova}
\begin{tabular}{@{}lrrr@{}}
\toprule
Kaynak & Kareler toplamı & sd & Kareler ortalaması \\
\midrule
Regresyon & 5312,182953 & 1 & 5312,182953 \\
Artık & 2957,725907 & 393 & 7,526020 \\
Toplam & 8269,908861 & 394 & --- \\
\bottomrule
\end{tabular}
\par\smallskip
{\footnotesize $F$, regresyon kareler ortalamasının artık kareler
ortalamasına oranıdır. Artık kareler ortalamasının karekökü 2,743359
model standart hatasını verir. Ara hesaplar yuvarlanmamış değerlerle
yapılır; tablolarda çoğu değer altı ondalık basamakla gösterilmiştir.\par}
\end{table}

\textbf{Grafik kontrolünün kapsamı.} Üç yazılımda da G1--G3
saçılımı, standartlaştırılmış artık--tahmin grafiği ve artık
histogramı paylaşılmıştır. Saçılımda pozitif doğrusal örüntü
görülür; G3=0 satırındaki noktalar ana örüntüden ayrılır.
Notlar tam sayı olduğundan noktalar üst üste binebilir; grafikteki
görünür simgelerden 395 kaydın veya sıfır notların sayısı okunmaz.

Artık--tahmin grafiklerinde ayrık bantlar ve yaklaşık $-4{,}24$'e
uzanan negatif artıklar vardır. Sıfır notlu gözlemlerde
artık $0-\widehat{G3}$ olduğundan aşağı eğimli bir sıra oluşur.
Histogramlarda sıfır çevresinde yoğunlaşma ve uzun negatif kuyruk
görülür. Bu görünüm normalliği ve sabit varyansı doğrulamış sayılmaz;
bantlar tek başına belirli bir varsayım ihlalinin kanıtı değildir.
Yazılımların histogram sınıf sınırları aynı olmayabilir; sütunların
farklı görünmesi sayısal sonuçların çeliştiği anlamına gelmez.
Aykırı görünen kayıtlar yalnız grafik nedeniyle silinmez.

\Needspace{12\baselineskip}
\begin{outputguidebox}
Tablo~\ref{tab:b13-dogrulanmis-katsayi} ile kurulan model
\[
\widehat{G3}=-1{,}652804+1{,}106256\,G1
\]
biçimindedir. G1'de bir puan daha yüksek değer, modelde ortalama
G3'ün yaklaşık 1,106 puan daha yüksek olmasıyla ilişkilidir;
G1'i bir puan artırmanın G3'ü bu miktarda artıracağı söylenmez.
G1'in gözlenen aralığı 3--19'dur; sabit terim G1=0'a karşılık
geldiğinden gözlenen aralığın dışındaki bir kesişimdir.

Tablo~\ref{tab:b13-dogrulanmis-uyum} içindeki $R^2=0{,}642351$,
bu örneklemde G3 değişkenliğinin yaklaşık yüzde 64,24'ünün modelle
açıklanan payıdır. Öğrencilerin yüzde 64'ünün doğru tahmin edildiği
veya yeni öğrencilerde aynı başarının elde edileceği anlamına gelmez.
G1=10 için 9,409757, modelin ortalama tahminidir; bireysel tahmin
aralığı değildir. Bu değer ara hesaplarda yuvarlanmamış katsayılarla
hesaplanmıştır.

Küçük $p$ değeri model varsayımlarını doğrulamaz. Negatif artık
kuyruğu klasik test ve aralıkların sınırlılıklarıyla birlikte
raporlanmalıdır. Farklı öğrencilerin bağımsızlığı grafiklerden
anlaşılamaz; okul, önceki koşullar ve diğer karıştırıcılar için
düzeltme yapılmamıştır. Notlar aynı eğitim sürecinin ilişkili
ölçümleridir. İki okulla sınırlı bu verideki uyum, nedensel etkiyi,
temsil gücünü veya yeni bir okulda tahmin başarısını kanıtlamaz.
\end{outputguidebox}

\Needspace{10\baselineskip}
\textbf{Raporlama örneği (doğrulanmış çıktılarla).}
``395 öğrencide birinci dönem ve yıl sonu notları arasında pozitif
doğrusal ilişki bulunmuştur ($r=0{,}801468$,
$p\approx9{,}00143\times10^{-90}$). Sabit terimli basit regresyonda
G1 eğimi 1,106256'dır (\%95 güven aralığı
$[1{,}024393;\ 1{,}188119]$); $R^2=0{,}642351$ ve
$F(1,393)=705{,}842247$ elde edilmiştir. Modelin artık standart
hatası 2,743359 puandır. Sıfır yıl sonu notu bulunan 38 kayıt
korunmuştur. Artık grafiklerinde negatif kuyruk ve uzak noktalar
görülmüş; normallik ve sabit varyans doğrulanmış kabul edilmemiştir.
Klasik test ve güven aralığı bu sınırlılıklarla yorumlanmıştır.
Okul ve diğer karıştırıcılar için düzeltme yapılmamış; bulgular
nedensel etki veya yeni öğrencilerde doğrulanmış tahmin başarısı
olarak sunulmamıştır.''

\textbf{Görev.} Standartlaştırılmış Beta ile B eğimini ayırt edin.
Yuvarlanmamış sonuçlarla $r^2=R^2$ ve $t^2=F$ eşitliklerini kontrol
edin. G1=10 için ortalama tahmini ile bireysel tahmin aralığının
neden farklı olduğunu açıklayın. Artıkların negatif kuyruğunu
raporlayın; sıfır notları yalnız model uyumunu artırmak için silmeyin.
Eğitim yılı başlamadan yapılacak bir tahminde henüz oluşmamış \texttt{g1}
notunu kullanmanın neden uygun olmayacağını açıklayın.